
\begin{tikzpicture}[use Hobby shortcut]

    % ============================================================
    % PARAMETRES GENERAUX
    % ============================================================

    \def\longueur{3.0}
    \def\dx{0.18}
    \def\dy{0.10}

    % ============================================================
    % TITRE
    % ============================================================

    \node[
        color=\colorOne,
        align=center
    ] at (5,10.5)
    {
        \fontsize{25pt}{30pt}\selectfont
        Pourquoi peut-on écrire
        $\displaystyle \mu=\rho\,S$ ?
    };


    % ============================================================
    % 1. CORDE NON HOMOGENE
    % ============================================================

    \begin{scope}[shift={(1.7,5.2)}]

        % Titre
        \node[
            color=\colorOne,
            align=center
        ] at (0,3.0)
        {
            \fontsize{16pt}{20pt}\selectfont
            \textbf{1. Section non homogène}
        };

        % --------------------------------------------------------
        % Direction de la corde
        % --------------------------------------------------------

        \draw[
            dashed,
            color=\colorThree,
            line width=0.2ex
        ]
        (-1.5,-1.0) -- (3.5,1.8);

        % --------------------------------------------------------
        % Sections successives avec taille variable
        % et densité variable
        % --------------------------------------------------------

        \foreach \i in {12,...,0}
        {

            % Position dans la direction diagonale
            \pgfmathsetmacro{\xshift}{\i*\dx}
            \pgfmathsetmacro{\yshift}{\i*\dy}

            % Taille légèrement irrégulière
            \pgfmathsetmacro{\s}{1.15-0.018*\i}

            % Variation de couleur :
            % cœur clair -> extérieur rouge
            \pgfmathsetmacro{\c}{10+6*\i}

            \begin{scope}[
                shift={(\xshift,\yshift)},
                scale=\s
            ]

                % Section irrégulière
                \draw[
                    closed,
                    fill=colorFive!\c!colorTwo,
                    fill opacity=0.90,
                    draw=colorOne,
                    line width=0.18ex
                ]
                (0,-1.00)
                .. (0.85,-1.10)
                .. (1.25,-0.40)
                .. (1.10,0.55)
                .. (0.45,1.00)
                .. (-0.35,0.85)
                .. (-1.10,0.40)
                .. (-1.25,-0.35)
                .. (-0.70,-0.95)
                -- cycle;

            \end{scope}
        }

        % --------------------------------------------------------
        % Indication rho(x,y)
        % --------------------------------------------------------

        \node[
            color=\colorSix,
            fill=white,
            fill opacity=0.85,
            text opacity=1,
            inner sep=4pt,
            align=center
        ]
        at (0,-2.0)
        {
            \fontsize{13pt}{16pt}\selectfont
            $\rho=\rho(x,y)$
        };

        % --------------------------------------------------------
        % Fleche vers formule
        % --------------------------------------------------------

        \draw[
            color=\colorFour,
            line width=0.35ex,
            ->,
            >=Stealth
        ]
        (0,-1.5)
        to[out=-90,in=180]
        (1.0,-2.8);

        % Formule
        \node[
            color=\colorFour,
            align=left
        ]
        at (2.2,-2.9)
        {
            \fontsize{12pt}{15pt}\selectfont
            $\displaystyle
            \mu
            =
            \iint_S
            \rho(x,y)\,dS$
        };

    \end{scope}


    % ============================================================
    % 2. CORDE HOMOGENE - SECTION NON CIRCULAIRE
    % ============================================================

    \begin{scope}[shift={(5.0,5.2)}]

        % Titre
        \node[
            color=\colorOne,
            align=center
        ] at (0,3.0)
        {
            \fontsize{16pt}{20pt}\selectfont
            \textbf{2. Section homogène}\\
            \textbf{non circulaire}
        };

        % --------------------------------------------------------
        % Axe de la corde
        % --------------------------------------------------------

        \draw[
            dashed,
            color=\colorThree,
            line width=0.2ex
        ]
        (-1.5,-1.0) -- (3.5,1.8);

        % --------------------------------------------------------
        % Empilement de sections elliptiques
        % --------------------------------------------------------

        \foreach \i in {15,...,0}
        {

            \pgfmathsetmacro{\xshift}{\i*\dx}
            \pgfmathsetmacro{\yshift}{\i*\dy}

            \begin{scope}[
                shift={(\xshift,\yshift)}
            ]

                % Ellipse représentant la section
                \draw[
                    fill=colorFour!20!colorTwo,
                    fill opacity=0.85,
                    draw=colorFour,
                    line width=0.25ex
                ]
                (0,0)
                ellipse
                [
                    x radius=1.15,
                    y radius=0.75
                ];

            \end{scope}
        }

        % --------------------------------------------------------
        % Demi-axes a et b
        % --------------------------------------------------------

        \draw[
            <->,
            >=Stealth,
            color=\colorOne,
            line width=0.25ex
        ]
        (-1.55,-0.75)
        --
        (-1.55,0.75)
        node[
            midway,
            left,
            color=\colorOne
        ]
        {
            \fontsize{13pt}{16pt}\selectfont
            $2a$
        };

        \draw[
            <->,
            >=Stealth,
            color=\colorOne,
            line width=0.25ex
        ]
        (-1.15,-1.10)
        --
        (1.15,-1.10)
        node[
            midway,
            below,
            color=\colorOne
        ]
        {
            \fontsize{13pt}{16pt}\selectfont
            $2b$
        };

        % --------------------------------------------------------
        % Homogeneite
        % --------------------------------------------------------

        \node[
            color=\colorFour,
            fill=white,
            fill opacity=0.85,
            text opacity=1,
            inner sep=4pt
        ]
        at (0,-2.0)
        {
            \fontsize{13pt}{16pt}\selectfont
            $\rho(x,y)=\rho=\text{constante}$
        };

        % --------------------------------------------------------
        % Formule
        % --------------------------------------------------------

        \draw[
            color=\colorFour,
            line width=0.35ex,
            ->,
            >=Stealth
        ]
        (0,-1.5)
        to[out=-90,in=180]
        (0.9,-2.8);

        \node[
            color=\colorFour,
            align=left
        ]
        at (2.0,-2.9)
        {
            \fontsize{12pt}{15pt}\selectfont
            $\displaystyle
            \mu
            =
            \rho\,S$
            \\[0.15cm]
            $\displaystyle
            S=\pi ab$
        };

    \end{scope}


    % ============================================================
    % 3. CORDE HOMOGENE - SECTION CIRCULAIRE
    % ============================================================

    \begin{scope}[shift={(8.7,5.2)}]

        % Titre
        \node[
            color=\colorOne,
            align=center
        ] at (0,3.0)
        {
            \fontsize{16pt}{20pt}\selectfont
            \textbf{3. Section homogène}\\
            \textbf{circulaire}
        };

        % --------------------------------------------------------
        % Axe de la corde
        % --------------------------------------------------------

        \draw[
            dashed,
            color=\colorThree,
            line width=0.2ex
        ]
        (-1.5,-1.0) -- (3.5,1.8);

        % --------------------------------------------------------
        % Empilement des sections circulaires
        % --------------------------------------------------------

        \foreach \i in {15,...,0}
        {

            \pgfmathsetmacro{\xshift}{\i*\dx}
            \pgfmathsetmacro{\yshift}{\i*\dy}

            \begin{scope}[
                shift={(\xshift,\yshift)}
            ]

                \draw[
                    fill=colorGold!30!colorTwo,
                    fill opacity=0.85,
                    draw=colorFour,
                    line width=0.25ex
                ]
                (0,0)
                circle
                [radius=1.0];

            \end{scope}
        }

        % --------------------------------------------------------
        % Diametre
        % --------------------------------------------------------

        \draw[
            <->,
            >=Stealth,
            color=\colorOne,
            line width=0.25ex
        ]
        (-1.45,-1.0)
        --
        (-1.45,1.0)
        node[
            midway,
            left
        ]
        {
            \fontsize{14pt}{17pt}\selectfont
            $d$
        };

        % --------------------------------------------------------
        % Homogeneite
        % --------------------------------------------------------

        \node[
            color=\colorFour,
            fill=white,
            fill opacity=0.85,
            text opacity=1,
            inner sep=4pt
        ]
        at (0,-2.0)
        {
            \fontsize{13pt}{16pt}\selectfont
            $\rho(x,y)=\rho=\text{constante}$
        };

        % --------------------------------------------------------
        % Formules finales
        % --------------------------------------------------------

        \draw[
            color=\colorFour,
            line width=0.35ex,
            ->,
            >=Stealth
        ]
        (0,-1.5)
        to[out=-90,in=180]
        (0.8,-2.8);

        \node[
            color=\colorFour,
            align=left
        ]
        at (2.0,-2.9)
        {
            \fontsize{12pt}{15pt}\selectfont
            $\displaystyle
            S
            =
            \frac{\pi d^2}{4}$
            \\[0.15cm]
            $\displaystyle
            \boxed{
            \mu
            =
            \rho
            \frac{\pi d^2}{4}
            }$
        };

    \end{scope}


    % ============================================================
    % FLECHES LOGIQUES ENTRE LES TROIS CAS
    % ============================================================

    \draw[
        ->,
        >=Stealth,
        color=\colorSix,
        line width=0.55ex
    ]
    (3.0,5.0)
    --
    (3.8,5.0);

    \node[
        color=\colorSix,
        align=center
    ]
    at (3.4,5.45)
    {
        \fontsize{10pt}{12pt}\selectfont
        matériau\\homogène
    };


    \draw[
        ->,
        >=Stealth,
        color=\colorSix,
        line width=0.55ex
    ]
    (6.6,5.0)
    --
    (7.5,5.0);

    \node[
        color=\colorSix,
        align=center
    ]
    at (7.05,5.45)
    {
        \fontsize{10pt}{12pt}\selectfont
        section\\circulaire
    };


    % ============================================================
    % CHAINE LOGIQUE FINALE
    % ============================================================

    \node[
        draw=\colorFour,
        line width=0.25ex,
        rounded corners=5pt,
        fill=colorFour!8!white,
        align=center,
        text width=14cm,
        inner sep=10pt,
        color=\colorOne
    ]
    at (5.0,0.5)
    {
        \fontsize{15pt}{20pt}\selectfont
        \[
        \mu
        =
        \iint_S
        \rho(x,y)\,dS
        \quad
        \xrightarrow[\text{matériau homogène}]{\rho=\text{constante}}
        \quad
        \mu=\rho S
        \quad
        \xrightarrow[\text{section circulaire}]{S=\pi d^2/4}
        \quad
        \boxed{
        \mu
        =
        \rho\frac{\pi d^2}{4}
        }
        \]
    };

\end{tikzpicture}


\begin{tikzpicture}

% ============================================================
% PARAMETRES
% ============================================================

% Direction de la longueur de la corde
\def\dx{0.18}
\def\dy{0.10}

% Style commun des trois cordes
\tikzset{
    corde contour/.style={
        line width=0.25ex,
        color=\colorOne
    }
}

% ============================================================
% TITRE
% ============================================================

\node[
    color=\colorOne,
    align=center
]
at (5,10.7)
{
    \fontsize{24pt}{28pt}\selectfont
    \textbf{De la masse volumique à la masse linéique}
};

% ============================================================
% ============================================================
% 1. CORDE NON HOMOGENE
% ============================================================
% ============================================================

\begin{scope}[shift={(1.8,6.0)}]

    % --------------------------------------------------------
    % Titre
    % --------------------------------------------------------

    \node[
        color=\colorOne,
        align=center
    ]
    at (0,2.7)
    {
        \fontsize{15pt}{18pt}\selectfont
        \textbf{1. Section non homogène}
    };


    % --------------------------------------------------------
    % Axe de la corde
    % --------------------------------------------------------

    \draw[
        dashed,
        color=\colorThree,
        line width=0.2ex
    ]
    (-1.7,-1.1)
    --
    (3.7,1.9);


    % --------------------------------------------------------
    % Couches successives
    %
    % Chaque couche possède :
    % - une position différente
    % - une taille différente
    % - une couleur différente
    %
    % Cela représente une densité non uniforme
    % --------------------------------------------------------

    \foreach \i in {18,...,0}
    {

        % Position de chaque section
        \pgfmathsetmacro{\xshift}{\i*\dx}
        \pgfmathsetmacro{\yshift}{\i*\dy}

        % Taille :
        % le bout visible est légèrement plus grand
        \pgfmathsetmacro{\rx}{1.15-0.008*\i}
        \pgfmathsetmacro{\ry}{0.90-0.006*\i}

        % Variation de couleur
        \pgfmathsetmacro{\mix}{10+4*\i}

        \begin{scope}[
            shift={(\xshift,\yshift)}
        ]

            % Section irrégulière
            %
            % Ici on n'utilise PAS Hobby.
            % On ferme simplement avec -- cycle.

            \path[
                fill=colorSix!\mix!colorTwo,
                fill opacity=0.90,
                corde contour
            ]
            (0,-\ry)
            .. controls
                (0.70*\rx,-1.10*\ry)
                and
                (1.20*\rx,-0.45*\ry)
            ..
            (1.05*\rx,0.40*\ry)

            .. controls
                (0.85*\rx,1.05*\ry)
                and
                (0.30*\rx,1.15*\ry)
            ..
            (-0.20*\rx,0.95*\ry)

            .. controls
                (-0.90*\rx,0.90*\ry)
                and
                (-1.25*\rx,0.30*\ry)
            ..
            (-1.10*\rx,-0.35*\ry)

            .. controls
                (-0.90*\rx,-0.95*\ry)
                and
                (-0.40*\rx,-1.10*\ry)
            ..
            (0,-\ry)

            -- cycle;

        \end{scope}
    }


    % --------------------------------------------------------
    % Densite variable
    % --------------------------------------------------------

    \node[
        color=\colorSix,
        fill=white,
        rounded corners=2pt,
        inner sep=5pt
    ]
    at (0,-1.8)
    {
        \fontsize{14pt}{16pt}\selectfont
        $\rho=\rho(x,y)$
    };


    % --------------------------------------------------------
    % Formule generale
    % --------------------------------------------------------

    \node[
        color=\colorOne,
        align=center
    ]
    at (0,-3.0)
    {
        \fontsize{13pt}{16pt}\selectfont
        $\displaystyle
        \mu
        =
        \iint_S
        \rho(x,y)\,dS$
    };

\end{scope}


% ============================================================
% ============================================================
% 2. CORDE HOMOGENE
% SECTION NON CIRCULAIRE
% ============================================================
% ============================================================

\begin{scope}[shift={(5.0,6.0)}]

    % --------------------------------------------------------
    % Titre
    % --------------------------------------------------------

    \node[
        color=\colorOne,
        align=center
    ]
    at (0,2.7)
    {
        \fontsize{15pt}{18pt}\selectfont
        \textbf{2. Section homogène}\\
        \textbf{non circulaire}
    };


    % --------------------------------------------------------
    % Axe de la corde
    % --------------------------------------------------------

    \draw[
        dashed,
        color=\colorThree,
        line width=0.2ex
    ]
    (-1.7,-1.1)
    --
    (3.7,1.9);


    % --------------------------------------------------------
    % Corde :
    % empilement d'ellipses identiques
    % --------------------------------------------------------

    \foreach \i in {18,...,0}
    {

        \pgfmathsetmacro{\xshift}{\i*\dx}
        \pgfmathsetmacro{\yshift}{\i*\dy}

        \begin{scope}[
            shift={(\xshift,\yshift)}
        ]

            \path[
                fill=colorFour!25!white,
                fill opacity=0.90,
                corde contour
            ]
            (0,0)
            ellipse[
                x radius=1.15,
                y radius=0.75
            ];

        \end{scope}
    }


    % --------------------------------------------------------
    % Axes de la section
    % --------------------------------------------------------

    \draw[
        <->,
        >=Stealth,
        color=\colorOne,
        line width=0.22ex
    ]
    (-1.55,-0.75)
    --
    (-1.55,0.75)

    node[
        midway,
        left
    ]
    {
        \fontsize{12pt}{14pt}\selectfont
        $2a$
    };


    \draw[
        <->,
        >=Stealth,
        color=\colorOne,
        line width=0.22ex
    ]
    (-1.15,-1.05)
    --
    (1.15,-1.05)

    node[
        midway,
        below
    ]
    {
        \fontsize{12pt}{14pt}\selectfont
        $2b$
    };


    % --------------------------------------------------------
    % Homogeneite
    % --------------------------------------------------------

    \node[
        color=\colorFour,
        fill=white,
        rounded corners=2pt,
        inner sep=5pt
    ]
    at (0,-1.8)
    {
        \fontsize{13pt}{15pt}\selectfont
        $\rho(x,y)=\rho=\mathrm{constante}$
    };


    % --------------------------------------------------------
    % Formules
    % --------------------------------------------------------

    \node[
        color=\colorOne,
        align=center
    ]
    at (0,-3.0)
    {
        \fontsize{13pt}{16pt}\selectfont
        $\displaystyle
        \mu=\rho S$

        \vspace{0.1cm}

        $\displaystyle
        S=\pi ab$
    };

\end{scope}


% ============================================================
% ============================================================
% 3. CORDE HOMOGENE
% SECTION CIRCULAIRE
% ============================================================
% ============================================================

\begin{scope}[shift={(8.5,6.0)}]

    % --------------------------------------------------------
    % Titre
    % --------------------------------------------------------

    \node[
        color=\colorOne,
        align=center
    ]
    at (0,2.7)
    {
        \fontsize{15pt}{18pt}\selectfont
        \textbf{3. Section homogène}\\
        \textbf{circulaire}
    };


    % --------------------------------------------------------
    % Axe de la corde
    % --------------------------------------------------------

    \draw[
        dashed,
        color=\colorThree,
        line width=0.2ex
    ]
    (-1.7,-1.1)
    --
    (3.7,1.9);


    % --------------------------------------------------------
    % Corde circulaire
    % --------------------------------------------------------

    \foreach \i in {18,...,0}
    {

        \pgfmathsetmacro{\xshift}{\i*\dx}
        \pgfmathsetmacro{\yshift}{\i*\dy}

        \begin{scope}[
            shift={(\xshift,\yshift)}
        ]

            \path[
                fill=colorGold!35!white,
                fill opacity=0.90,
                corde contour
            ]
            (0,0)
            circle[
                radius=1
            ];

        \end{scope}
    }


    % --------------------------------------------------------
    % Diametre d
    % --------------------------------------------------------

    \draw[
        <->,
        >=Stealth,
        color=\colorOne,
        line width=0.22ex
    ]
    (-1.40,-1.0)
    --
    (-1.40,1.0)

    node[
        midway,
        left
    ]
    {
        \fontsize{15pt}{18pt}\selectfont
        $d$
    };


    % --------------------------------------------------------
    % Centre
    % --------------------------------------------------------

    \fill[
        color=\colorSix
    ]
    (0,0)
    circle[
        radius=0.06
    ];


    % --------------------------------------------------------
    % Rayon
    % --------------------------------------------------------

    \draw[
        color=\colorSix,
        line width=0.25ex,
        ->
    ]
    (0,0)
    --
    (0.75,0.65)

    node[
        midway,
        above,
        color=\colorSix
    ]
    {
        \fontsize{11pt}{13pt}\selectfont
        $R$
    };


    % --------------------------------------------------------
    % Homogeneite
    % --------------------------------------------------------

    \node[
        color=\colorFour,
        fill=white,
        rounded corners=2pt,
        inner sep=5pt
    ]
    at (0,-1.8)
    {
        \fontsize{13pt}{15pt}\selectfont
        $\rho(x,y)=\rho=\mathrm{constante}$
    };


    % --------------------------------------------------------
    % Resultat
    % --------------------------------------------------------

    \node[
        color=\colorOne,
        align=center
    ]
    at (0,-3.0)
    {
        \fontsize{13pt}{16pt}\selectfont
        $\displaystyle
        S
        =
        \pi R^2
        =
        \frac{\pi d^2}{4}$
    };

\end{scope}


% ============================================================
% FLECHES DE PROGRESSION
% ============================================================

\draw[
    ->,
    >=Stealth,
    line width=0.5ex,
    color=\colorSix
]
(3.2,6.2)
--
(3.75,6.2);

\node[
    color=\colorSix,
    align=center
]
at (3.48,6.65)
{
    \fontsize{9pt}{11pt}\selectfont
    matériau\\homogène
};


\draw[
    ->,
    >=Stealth,
    line width=0.5ex,
    color=\colorSix
]
(6.75,6.2)
--
(7.30,6.2);

\node[
    color=\colorSix,
    align=center
]
at (7.02,6.65)
{
    \fontsize{9pt}{11pt}\selectfont
    section\\circulaire
};


% ============================================================
% RESULTAT FINAL
% ============================================================

\node[
    draw=\colorFour,
    line width=0.3ex,
    rounded corners=8pt,
    fill=colorFour!8!white,
    text width=15cm,
    align=center,
    inner sep=8pt,
    color=\colorOne
]
at (5,1.0)
{
    \fontsize{16pt}{20pt}\selectfont
    \textbf{Cas de la corde en nylon de l'exercice}

    \vspace{0.25cm}

    \[
        \boxed{
            \mu
            =
            \rho S
            =
            \rho\frac{\pi d^2}{4}
        }
    \]
};


% ============================================================
% LEGENDE
% ============================================================

\node[
    color=\colorThree,
    align=center
]
at (5,-0.3)
{
    \fontsize{11pt}{14pt}\selectfont
    L'homogénéité permet d'écrire
    $\mu=\rho S$ ;
    la forme circulaire permet de calculer
    $S=\pi d^2/4$.
};

\end{tikzpicture}







\begin{tikzpicture}[use Hobby shortcut]

	%\fill[fill=\colorFour , decoration={random steps,segment length=2mm} , decorate , closed] (-1,-1) -- (-1,11) --(11,11) -- (11,-1)-- cycle ;
	%\fill[fill=\colorFour ] (0,0) -- (0,10)  -- (10,0) ;

	\begin{scope}[shift={(4.5,5)} , scale = 1.0]
		%\draw[fill= black ] (-5,-5) rectangle + (10 , 10) ;
		\newcommand{\curve}{(-2,-1) .. (1,-2) .. (2,2) .. (0,1)}

        \pgfmathsetmacro{\dx}{0.1}
        \pgfmathsetmacro{\dy}{0.1}
		\foreach \i in {0,...,5} { % 50 couches
				\pgfmathsetmacro{\s}{0.2 + 0.5*\i}
				\pgfmathsetmacro{\c}{2*\i}   % mélange de 0% à 98%
                \pgfmathsetmacro{\xshift}{1.0*\i}
                \pgfmathsetmacro{\yshift}{-0.5*\i}
                %\draw[closed, scale=\s, fill=colorSix!\c!colorTwo, draw=colorSix!\c!colorTwo , shift={(\xshift,\yshift)}] \curve -- cycle;
	
				\draw[closed, scale=\s, fill=none, draw=colorSix , shift={(\xshift,\yshift)}] \curve -- cycle;
			}

		\begin{scope}[shift={(-4.5,1)} , scale = 1.0]
			\draw[fill=\colorFour , opacity = 0.5 , draw = none  ]
			(1.9,0) rectangle (9.85, 0.5 )
			;
			\draw[line width=0.5ex, color=\colorFour , dashed]
			(2,0) -- (9.75,0)
			(2,0.5) -- (9.75,0.5)
			;
			\draw[color=\colorFour]
			(10,0) edge[<->, line width=0.1ex , ultra thick] node[pos= 0.5 , right]{\fontsize{25pt}{16pt}\selectfont $\delta \theta$} (10,0.5)
			;
			\begin{scope}%[overlay]
				\node[] (point) at (4,0.1) {};
				\draw[  opacity = 1 ]
				(point) edge[line width=0.4ex ,color = \colorFour ,arrows = {Computer Modern Rightarrow[sharp ,length=0.2cm]-} , out = 90 , in = 180  ]node[pos=1, right , font = \fontsize{12pt}{16pt}\selectfont]{\color{colorFour} \fontsize{20}{10}\selectfont \bfseries \begin{tabular}{c} $\displaystyle \int dx \rho ( x , \theta ; t ) \, \delta \theta = \Pi(\theta) \, \delta \theta  $  \\[0.2cm] \textbf{est conservée} \end{tabular} }($(point) + (2.0cm,2.0cm)$)
				;
			\end{scope}
		\end{scope}

	\end{scope}

	\node[draw = none , fill=none, minimum width=3cm, minimum height=2cm, align=center, \colorTwo]
	at (5,4) {\fontsize{20pt}{16pt}\selectfont $\rho(x , \theta ; t )$};
	%\node[draw = none , fill=none, minimum width=3cm, minimum height=2cm, align=center, \colorTwo] 
	%    at (8,1.5) {\fontsize{20pt}{16pt}\selectfont $\nu = 0 $};

	% \node[draw = none , fill=none, minimum width=3cm, minimum height=2cm, align=center , \colorTwo , rotate = 40] 
	%     at (5,4.5) {\fontsize{20pt}{16pt}\selectfont $\nu = 1$};

	%    \draw[ ]
	%    		%(8,6.5) edge[line width=0.3ex, color=\colorThree , dashed ] node [pos = 0 , below ,\colorThree ]{\fontsize{10pt}{16pt}\selectfont $x_2$} (8,9)
	%    		%(8,7.25) edge[line width=0.3ex, color=\colorThree , ->] node [pos = 1 , below ,\colorThree ]{\fontsize{10pt}{16pt}\selectfont $v^{\mathrm{eff}}_{[\rho ]}(\theta_+)$}(9,7.25) 
	%    		(3,1.5) edge[line width=0.5ex, color=\colorGold!50!\colorThree , dashed ] node [pos = 0 , below ,\colorGold!50!\colorThree]{\fontsize{15pt}{16pt}\selectfont $x_1$}(3,6)
	%    		(3,5.4) edge[line width=0.5ex, color=\colorGold!50!\colorThree , ->] node [pos = 1 , below ,\colorGold!50!\colorThree]{\fontsize{15pt}{16pt}\selectfont $v^{\mathrm{eff}}_{[\rho  ]}(\theta_-)$}(2.0,5.4) 
	%    ;

	% Axe horizontal (gamma)
	\draw[ line width=0.8ex, color=\colorOne , ->]
	(-0.5,0) -- (10.5,0)
	node[shift = {(0,0)} , pos=0.5, below] {\fontsize{40pt}{19pt}\selectfont $x$};

	% Axe vertical (t)
	\draw[ line width=0.8ex, color=\colorOne , ->]
	(0,-0.5) -- (0,10.5)
	node[ shift = {(0,0)} , pos=0.5, above, rotate=90] {\fontsize{40pt}{19pt}\selectfont $\theta$};

	\drawgrid{-2}{11}{-2}{2}{1}{1};
\end{tikzpicture}


\begin{tikzpicture}[3d view={120}{15}, scale=1.5]

    % Couleurs personnalisées pour donner du relief à la corde
    \definecolor{ropeLight}{HTML}{D2B48C}
    \definecolor{ropeDark}{HTML}{8B5A2B}
    \definecolor{ropeCut}{HTML}{EEDC82}

    % ==========================================
    % 1. CORPS DE LA CORDE (Cylindre principal)
    % ==========================================
    % Face arrière (fond de la corde)
    \draw[fill=ropeDark, draw=ropeDark!80!black] (0,0,0) circle (1);
    
    % Longueur de la corde (lignes de fuite)
    \draw[left color=ropeDark, right color=ropeLight, draw=ropeDark!50] 
        (0,1,0) -- (0,1,4) 
        arc[start angle=90, end angle=270, radius=1] 
        -- (0,-1,0) 
        arc[start angle=270, end angle=90, radius=1];

    % ==========================================
    % 2. COUPE 3D (Plan avant à z=4)
    % ==========================================
    \begin{scope}[canvas is xy plane at z=4]
        
        % Fond de la section coupée
        \draw[fill=ropeDark!60, draw=black, thick] (0,0) circle (1);
        
        % Toron 1 (Haut)
        \draw[fill=ropeCut, draw=ropeDark] (0, 0.4) circle (0.45);
        \draw[ropeDark!70, thin] (0, 0.4) circle (0.35);
        \draw[ropeDark!70, thin] (0, 0.4) circle (0.2);
        
        % Toron 2 (Bas Gauche)
        \draw[fill=ropeCut, draw=ropeDark] (-0.35, -0.2) circle (0.45);
        \draw[ropeDark!70, thin] (-0.35, -0.2) circle (0.35);
        \draw[ropeDark!70, thin] (-0.35, -0.2) circle (0.2);
        
        % Toron 3 (Bas Droite)
        \draw[fill=ropeCut, draw=ropeDark] (0.35, -0.2) circle (0.45);
        \draw[ropeDark!70, thin] (0.35, -0.2) circle (0.35);
        \draw[ropeDark!70, thin] (0.35, -0.2) circle (0.2);
        
        % Petit toron central de remplissage (optionnel)
        \draw[fill=ropeDark, draw=ropeDark] (0,0) circle (0.15);
        
    \end{scope}

    % ==========================================
    % 3. EFFET DE SERRAGE (Hachures extérieures)
    % ==========================================
    % Ajoute des lignes courbes sur le cylindre pour simuler la torsion de la corde
    \foreach \z in {0.5, 1.2, 1.9, 2.6, 3.3} {
        \draw[ropeDark!40, thick] (0,1,\z) to[out=-60, in=120] (0,-1,\z+0.4);
    }

\end{tikzpicture}